\documentclass[11pt,letterpaper]{article}
\usepackage[margin=1in]{geometry}
\usepackage{amsmath,amssymb,amsthm}
\usepackage[T1]{fontenc}
\pdfmapfile{+lm.map}
\pdfmapfile{+cm.map}
\pdfmapfile{+symbols.map}
\usepackage{lmodern,microtype}
\usepackage[hidelinks]{hyperref}
\hypersetup{pdftitle={Logarithmic deficit of 120 avoiding ascent sequences},pdfsubject={Two sided stretched bounds and inverse order for OEIS A202061}}
\allowdisplaybreaks[2]
\newtheorem{theorem}{Theorem}[section]
\newtheorem{lemma}[theorem]{Lemma}
\newtheorem{proposition}[theorem]{Proposition}
\newtheorem{corollary}[theorem]{Corollary}
\theoremstyle{remark}\newtheorem{remark}[theorem]{Remark}
\newcommand{\asc}{\operatorname{asc}}
\newcommand{\Q}{\mathbb Q}
\newcommand{\C}{\mathbb C}
\newcommand{\N}{\mathbb N}
\newcommand{\one}{\mathbf 1}
\newcommand{\st}{\star}
\title{Logarithmic deficit of 120 avoiding ascent sequences}
\author{A proof and reproducible research report}
\date{2 October 2026}
\begin{document}
\maketitle
\begin{abstract}
For the number $a_n$ of classical 120-avoiding ascent sequences, we derive an exact positive height-walk representation and prove
\[
 n\log\mu-\log a_n=\Theta\!\left(n^{1/3}(\log n)^{2/3}\right),
 \qquad \mu=7.295896943239772\ldots.
\]
This rules out the previously estimated pure $n^{3/8}$ stretched correction. The proof uses a positive binomial block formula, a critical jump tail, a finite-height radius bound, and a legal all-length staircase construction. We also obtain the order of the inverse correction and prove that the generating function is not D-finite. No exact stretched constant, asymptotic equivalent, or all-orders transseries is claimed. Exact-arithmetic certificates and independent finite diagnostics accompany the argument.
\end{abstract}

\section{Results and scope}
An ascent sequence is a word $w_1\cdots w_n$ of nonnegative integers with $w_1=0$ and
\[
 w_i\le1+\asc(w_1\cdots w_{i-1})\quad(i\ge2),\qquad
 \asc(w)=\#\{i:w_i<w_{i+1}\}.
\]
It avoids the classical pattern 120 if there are no $i<j<k$ with $w_k<w_i<w_j$. Let $a_n$ count these words, including $a_0=1$, and let $A(x)=\sum_{n\ge0}a_nx^n$. This is OEIS A202061 \cite{OEIS}.

\begin{theorem}\label{thm:main}
Let $\mu$ be the largest root of
\[
 \mu^3-8\mu^2+5\mu+1=0,
 \qquad \mu=7.29589694323977237457\ldots,
\]
and put $\Phi(n)=n^{1/3}(\log n)^{2/3}$. There exist $c,C>0$ and $n_0$ such that, for every integer $n\ge n_0$,
\begin{equation}\label{eq:main}
 \mu^n e^{-C\Phi(n)}\le a_n\le C\mu^n e^{-c\Phi(n)}.
\end{equation}
In particular,
\[
 n\log\mu-\log a_n=\Theta(\Phi(n)),\qquad
 \frac{\log(n\log\mu-\log a_n)}{\log n}\longrightarrow\frac13.
\]
The radius of convergence is $\rho=1/\mu$.
\end{theorem}

\begin{corollary}\label{cor:consequences}
The following consequences hold.
\begin{enumerate}
\item No equivalent $C_0\mu^n\exp(-d n^\gamma)n^g$, with $C_0,d>0$ and fixed $g\in\mathbb R$, is possible for $\gamma>1/3$ or for $\gamma=1/3$.
\item For $N(Y)=\min\{n\ge0:a_n\ge Y\}$ and $Y\to\infty$,
\[
 N(Y)=\frac{\log Y}{\log\mu}
 +\Theta\!\left((\log Y)^{1/3}(\log\log Y)^{2/3}\right).
\]
\item $A(x)$ is not D-finite over $\C(x)$; equivalently, $(a_n)$ is not P-recursive over $\C[n]$.
\end{enumerate}
\end{corollary}

Conway, Conway, Elvey Price and Guttmann \cite{Conway} proved equality of the exponential growth rates of the 120- and 201-avoiding classes. Their numerical analysis estimated a pure stretched power $\sigma=3/8$, with other fitted parameters depending on that choice. Their final paper says the later 500-term computation reinforced the estimate; it does not prove that stretched power. Theorem~\ref{thm:main} determines a different asymptotic order. It does not supply the leading constant in the deficit, prove convergence of the normalized deficit, or furnish a full expansion. The proof of its exponential rate is independent of any generating-function theorem for A202062.

The original derivation and the additional independent audit are included in the accompanying package. They are research verification, not journal peer review or formal proof-assistant certification. A scoped source check does not establish comprehensive publication priority; in particular the gap reduction may overlap an unpublished polynomial-time enumeration mentioned in \cite{Conway}.

\section{An exact normalized recurrence}
For an avoiding prefix, let $v$ be the largest lower endpoint of an already occurring increasing pair, or $0$ if there is none. A future letter smaller than $v$ would complete a 120 pattern, and every letter at least $v$ obeying the ascent bound is permitted. Delete the seen values smaller than $v$ and translate the surviving alphabet by $-v$. The normalized state is $(a,l,S)$, where $a$ is the number of previous ascents minus $v$, $l$ is the translated last letter, and $S$ is the set of seen translated surviving values. Write
\[
 S=\{0=s_0<s_1<\cdots<s_k\},\quad m=s_k,\quad
 h=a+1-m,\quad d_i=s_i-s_{i-1}.
\]
Every ascent sequence satisfies $\max w\le\asc(w)$, so $h\ge1$.

Appending $i\in\{0,\ldots,a+1\}$ sets
\[
 p=\max\bigl(\{s\in S:s<i\}\cup\{0\}\bigr)
\]
and gives the exact child
\begin{equation}\label{eq:state}
 (a',l',S')=
 \left(a+\one_{l<i}-p,\ i-p,\ \{s-p:s\in S,\ s\ge p\}\cup\{i-p\}\right).
\end{equation}
Indeed $p$ is the greatest newly forbidden lower threshold, and the insertion of $i-p$ records the new letter. Thus this recurrence follows directly from the class definition. In every normalized state, the last letter is either $0$ or the least positive member of $S$: if $i>0$, there was no old seen value strictly between $p$ and $i$.

Let $E_h(R)$ be the suffix generating function when the last letter is $0$ and the list of positive gaps is $R$. Let $G_h(d,R)$ be the suffix generating function when the gap list is $(d,R)$ and the last letter is $d$. All empty suffixes have weight $1$. Define
\[
 D_h(d,R)=\sum_{q=1}^dG_h(q,(d-q,R)),\quad e_h=E_h(\varnothing),\quad
 N_h=\sum_{q=1}^hG_{h+1-q}(q,\varnothing),
\]
where a zero gap is omitted. Write $(Tf)_h=f_{h+1}$.

\begin{lemma}\label{lem:operator}
Put $c=x/(1-x)$ and $\mathcal B=x+x^2T/(1-x)$. There are formal series $R_d(T),P_d(T),Q_d(T)$ with nonnegative integer coefficients in $x,T$ such that
\begin{align}
 D_h(d,R)&=R_d(T)E_h(R),\notag\\
 E_h((d,R))&=P_d(T)E_h(R),\qquad
 G_h(d,R)=Q_d(T)E_h(R),\label{eq:op}
\end{align}
where
\begin{align}
 P_d&=1+cTR_d,\qquad Q_d=1+\mathcal B R_d,\notag\\
 (1-\mathcal B)R_d&=1+\sum_{q=1}^{d-1}Q_qP_{d-q}.
 \label{eq:oprec}
\end{align}
\end{lemma}
\begin{proof}
Classify the next letter as zero, in an old gap, or a new maximum. For $R=(d_1,\ldots,d_k)$ this gives
\begin{equation}\label{eq:E}
 (1-x)E_h(R)=1+xN_h+x\sum_{j=1}^kD_{h+1}(d_j,(d_{j+1},\ldots,d_k)).
\end{equation}
From last letter $d$, selecting in the first gap gives no ascent; selecting in a later gap does. Subtraction from \eqref{eq:E} therefore gives
\[
 G_h(d,R)=E_h((d,R))+x(1-T)D_h(d,R).
\]
Subtracting \eqref{eq:E} for $R$ from its version for $(d,R)$ also gives
\[
 E_h((d,R))=E_h(R)+cTD_h(d,R).
\]
Together these imply $G_h(d,R)=E_h(R)+\mathcal B D_h(d,R)$. Induct on $d$, splitting the definition of $D$ into $q=d$ and $q<d$. For $q<d$, both shorter-gap identities apply to the arbitrary tail. All resulting operators commute, being series in the same shift $T$. The $q=d$ term gives the factor $1-\mathcal B$, whose inverse exists $x$-adically and has nonnegative coefficients. This proves the identities, existence, uniqueness and positivity.
\end{proof}

Use an ordinary variable $t$ for the shift and set $Q(x;z,t)=\sum_{q\ge1}Q_q(x,t)z^q$, with similar notation for $P,R$. With $b=1-x$ and $\mathcal B=x+x^2t/b$, summing \eqref{eq:oprec} yields
\[
 R=Q(1+P),\qquad Q=\frac{z}{1-z}+\mathcal B R,
 \qquad P=\frac{z}{1-z}+\frac{xt}{b}R.
\]
Elimination gives the exact jump equation
\begin{equation}\label{eq:Qquad}
 xt(1-z)Q^2+[b(z-b)+xt(x-z)]Q+bz=0,
\end{equation}
with $Q(x;0,t)=0$ and $Q(0;z,t)=z/(1-z)$. In particular,
\begin{equation}\label{eq:Q1}
 Q_1(x,t)=\frac1{1-x-x^2t/(1-x)}.
\end{equation}

\section{Positive height walks and block weights}
The empty-gap case of \eqref{eq:E} is
\begin{equation}\label{eq:walk}
 e_h=\frac1{1-x}+\frac{x}{1-x}\sum_{q=1}^h(Q_q(x,T)e)_{h+1-q},
 \qquad A(x)=1+xe_1(x).
\end{equation}
Interpret a term indexed by $q,r$ as a macro-step that first drops by $q-1$, then rises by $r$, has weight
\[
 \frac{x}{1-x}[t^r]Q_q(x,t),
\]
and changes height from $h$ to $h+1-q+r$. Its legality condition is exactly $1\le q\le h$. The terminal weight is $1/(1-x)$ at every height. This is a positive formal representation of the original sequence; no additional word-level bijection is required for coefficient selection. Every nonterminal step has positive $x$-degree, and only finitely many shifts occur at any fixed degree. Hence the formal Neumann series is well defined and uniquely solves \eqref{eq:walk}.

Put $\beta=xt/b$ and $\eta=x^2t/b^2$. Equation~\eqref{eq:Qquad} is equivalently
\begin{equation}\label{eq:W}
 Q=\frac{z+xW}{1-z},\qquad
 W=\frac{z}{b}\frac{(1+\beta)(1+W)}{1-\eta(1+W)},\qquad [z^0]W=0.
\end{equation}
Lagrange inversion gives, for $j\ge1$,
\begin{equation}\label{eq:lagrange}
 [z^j]W=\frac{(1+\beta)^j}{b^j}
 \sum_{k\ge0}\mathcal C(j,k)\eta^k,
 \quad \mathcal C(j,k)=\frac1j\binom{j+k-1}{k}\binom{j+k}{k+1}>0.
\end{equation}
For $\ell,r\ge0$ and $q\ge1$, write $B(\ell,q,r)=[x^\ell t^r]Q_q(x,t)$. Expanding \eqref{eq:lagrange} yields
\begin{equation}\label{eq:blockexact}
 B(\ell,q,r)=\one_{\ell=0,r=0}
 +\sum_{j=1}^q\ \sum_{k=0}^r
 \mathcal C(j,k)\binom{j}{r-k}
 \binom{j+\ell-2}{\ell-r-k-1}.
\end{equation}
The summation includes only $0\le r-k\le j$ and $\ell\ge r+k+1$. In particular, keeping $j=q$ and one admissible $k$ gives
\begin{equation}\label{eq:blocklower}
 B(\ell,q,r)\ge\mathcal C(q,k)\binom{q}{r-k}
 \binom{q+\ell-2}{\ell-r-k-1}.
\end{equation}

\section{An exact critical point}
Let $\alpha$ be the unique positive root of $7\alpha^3+14\alpha^2-7\alpha-1=0$ and define
\begin{equation}\label{eq:constants}
 \kappa=\frac{-21\alpha^2+17\alpha+5}{29},\qquad
 \mu=\frac{1+\alpha}{1-\alpha-\kappa},\qquad
 z_\st=\frac{\kappa(1-\alpha-\kappa)}{(\alpha-\kappa)(2\alpha+\kappa)}.
\end{equation}
The following bounds and identities are certified by rational interval arithmetic and polynomial reduction in the supplied program:
\begin{align*}
 \alpha&=0.5100033746178888\ldots,& \kappa&=0.2830305201361994\ldots,\\
 z_\st&=0.1980622641951617\ldots,& z_\st^3-5z_\st^2+6z_\st-1&=0,\\
 0&<\kappa<\alpha,& \alpha+\kappa&<1.
\end{align*}
The value of $\mu$ satisfies the cubic and interval in Theorem~\ref{thm:main}, which identify its largest real root. Descartes' rule gives uniqueness of the positive $\alpha$ root, isolated in
\[
 \frac{51000337}{10^8}<\alpha<\frac{51000339}{10^8}.
\]

For $L(v)=v\log v$, let
\begin{align*}
 S(a,d,u)={}&2L(a+u)-L(a)-2L(u)-L(d-u)-L(a-d+u)\\
 &+L(1+a)-L(1-d-u)-L(a+d+u).
\end{align*}
This is the factorial entropy of \eqref{eq:blocklower} divided by $\ell$, for $a=q/\ell$, $d=r/\ell$ and $u=k/\ell$. At $(\alpha,\alpha,\kappa)$,
\begin{equation}\label{eq:entropy}
 S_u=0,\qquad S_a+S_d=0,\qquad S_d=\log z_\st,\qquad S=\log\mu.
\end{equation}
For clarity, the two stationary identities reduce to
\begin{align*}
 (\alpha+\kappa)^2(\alpha-\kappa)(1-\alpha-\kappa)
 &=\kappa^3(2\alpha+\kappa),\\
 (\alpha+\kappa)^2(1+\alpha)(1-\alpha-\kappa)
 &=\alpha(\alpha-\kappa)(2\alpha+\kappa)^2.
\end{align*}
The displayed expression for $z_\st$ gives $S_d$. Euler's identity applied to the homogeneous factorial entropy then gives $S=\log((1+\alpha)/(1-\alpha-\kappa))$. These are finite algebraic verifications; global maximality of $S$ is not needed.

Set $\rho=1/\mu$, $b_\st=1-\rho$, $t_\st=1/z_\st$, and
\begin{align}\label{eq:critical}
 \beta_\st&=\frac{\rho}{b_\st z_\st},&
 \eta_\st&=\frac{\rho^2}{b_\st^2z_\st},&
 Q_\st&=\frac{b_\st z_\st}{b_\st z_\st+\rho},\\
 W_\st&=\frac{(1-z_\st)Q_\st-z_\st}{\rho},&
 m_\st&=\frac{\rho}{b_\st}t_\st Q_\st
 =\frac{\rho}{b_\st z_\st+\rho}.\notag
\end{align}
The exact certificate proves the fixed-point identity \eqref{eq:W} at these values, and
\begin{equation}\label{eq:margins}
 W_\st=1+\beta_\st>0,\quad
 \eta_\st(2+\beta_\st)^2=1,\quad
 0<\eta_\st(1+W_\st)<\tfrac12,\quad
 \tfrac25<m_\st<\tfrac12.
\end{equation}
Numerically, $W_\st=1.8019377358\ldots$ and $m_\st=0.4450418679\ldots$.

\begin{lemma}\label{lem:criticalbound}
The jump series converges at $(\rho,z_\st,t_\st)$, with $Q(\rho;z_\st,t_\st)\le Q_\st$. Moreover,
\begin{equation}\label{eq:ehcritical}
 e_h(\rho)\le\frac{t_\st^h}{b_\st(1-m_\st)},\qquad
 A(\rho)\le1+\frac{\rho t_\st}{b_\st(1-m_\st)}<\infty.
\end{equation}
\end{lemma}
\begin{proof}
Expand the denominator on the right of \eqref{eq:W} geometrically. Iteration starting from $0$ is nonnegative and bounded by the fixed-point supersolution $W_\st$, by \eqref{eq:margins}. Monotone convergence proves the assertion about $Q$. For the height weight $t_\st^h$, the sum of an unrestricted transition row divided by its starting weight is at most
\[
 \frac{\rho}{b_\st}\sum_{q\ge1,r\ge0}
 [t^r]Q_q(\rho,t)\,t_\st^{1-q+r}
 =\frac{\rho}{b_\st}t_\st Q(\rho;z_\st,t_\st)\le m_\st<1.
\]
Each legal row is a subsum. The terminal weight $1/b_\st$ is at most $t_\st^h/b_\st$, so the nonnegative Neumann series is bounded by a geometric series. This proves \eqref{eq:ehcritical}. No spectral-radius identification from a discriminant is being assumed.
\end{proof}

\section{The all length lower bound}
\begin{lemma}\label{lem:uniform}
For each fixed $D>0$ there exist $c_D,C_D>0$ and $\ell_D$ such that the conditions
\[
 \ell\ge\ell_D,\quad |q-\alpha\ell|\le2,\quad
 |k-\kappa\ell|\le1,\quad
 |r-q|\le D\sqrt{\ell\log\ell}
\]
imply admissibility in \eqref{eq:blocklower} and
\begin{equation}\label{eq:uniform}
 B(\ell,q,r)\ge c_D\ell^{-C_D}\mu^\ell z_\st^{r-q}.
\end{equation}
If $|r-q|\le D\sqrt\ell$, one can use the more precise factor $c_D\ell^{-3}$.
\end{lemma}
\begin{proof}
All normalized factorial arguments lie in one compact subset of the positive interior, for all sufficiently large $\ell$. The exact identities
\begin{align*}
 \mathcal C(q,k)&=\frac{q}{(q+k)(k+1)}\binom{q+k}{k}^{\!2},\\
 \binom{q+\ell-2}{\ell-r-k-1}
 &=\frac{(\ell-r-k)(q+r+k)}{(q+\ell)(q+\ell-1)}
 \binom{q+\ell}{\ell-r-k}
\end{align*}
and uniform Stirling bounds give a lower bound
$c\ell^{-3}\exp(\ell S(q/\ell,r/\ell,k/\ell))$. Taylor's theorem and \eqref{eq:entropy} yield
\[
 \ell S=\ell\log\mu+(r-q)\log z_\st+O_D(\log\ell).
\]
Indeed the first differential is exactly $(r-q)\log z_\st$, and the quadratic remainder is $O((r-q)^2/\ell+1)$. This proves \eqref{eq:uniform}; with square-root perturbations the remainder is bounded.
\end{proof}

Choose a sufficiently large fixed integer $J$. For $j\ge J$, define
\[
 h_j=\lfloor j^2\log j\rfloor,\quad q_j=\lfloor h_j/2\rfloor,
 \quad \ell_j=\lfloor q_j/\alpha\rfloor,\quad
 k_j=\text{a nearest integer to }\kappa\ell_j.
\]
An ascent step from $h_j$ to $h_{j+1}$ uses
\[
 r_j^+=q_j-1+(h_{j+1}-h_j),
\]
and a descent step from $h_j$ to $h_{j-1}$ uses
\[
 r_j^-=q_j-1+(h_{j-1}-h_j).
\]
Both preserve the exact increment $1-q+r$. As $\ell_j=\Theta(j^2\log j)$ and $h_{j+1}-h_j=O(j\log j)$, every such step satisfies Lemma~\ref{lem:uniform} with the same $D$. Its down parameter obeys $q_j\le h_j$, so it is legal. Retain just the leading $x$ in the macro factor $x/(1-x)$; then its length is $\ell_j+1$ and multiplicity at least $B(\ell_j,q_j,r_j^{\pm})$.

Seed the walk from height $1$ to $h_J$ using $q=1,r=h_J-1$. Equation~\eqref{eq:Q1} gives a coefficient $1$ at its leading degree $\ell=2(h_J-1)$, so the macro-step has length $2h_J-1$. The final return from $h_J$ to $1$ uses $q=h_J,r=0,\ell=0$, of length $1$ and multiplicity $1$.

Ascend from $h_J$ to $h_K$ and descend to $h_J$, then add the seeds and the original leading letter. The total length is
\begin{align}\label{eq:LK}
 L_K&=2h_J+1+\sum_{j=J}^{K-1}(\ell_j+1)
                +\sum_{j=J+1}^{K}(\ell_j+1)\notag\\
 &=\frac{K^3\log K}{3\alpha}+O(K^3),
\end{align}
and
\[
 L_{K+1}-L_K=\ell_K+\ell_{K+1}+2=\Theta(h_K).
\]
The total increment of the large blocks is zero. If there are $M$ such blocks, then $\sum(r-q)=-M$, so their height-tilt factors multiply to $z_\st^{-M}$. There is no uncancelled exponentially expensive endpoint factor.

To cover every sufficiently large integer $n$, choose maximal $K$ with $L_K\le n$ and let $R=n-L_K$. A bounded residual is absorbed by the terminating geometric series. Otherwise split $R=m_1+m_2$ into two nearly equal integers, both above a fixed threshold, and insert two self-loops at height $h_K$. For length $m$ choose
\[
 \ell=m-1,\qquad q=\text{a nearest integer to }\alpha\ell,
 \qquad r=q-1,\qquad k=\text{a nearest integer to }\kappa\ell.
\]
The loop has exactly zero increment and meets Lemma~\ref{lem:uniform}. Furthermore,
\[
 q\le\frac{\alpha R}{2}+O(1)
 \le\frac{h_K+h_{K+1}}4+O(1)<h_K
\]
for all sufficiently large $K$, proving legality.

There are $O(K)$ large blocks and at most two repair loops, each with $\ell=O(K^2\log K)$. Multiply \eqref{eq:uniform}. The extra factor $\mu^{-1}$ for each macro-step's leading $x$, and the fixed seed lengths, cost at most $\exp(O(K))$. Polynomial factors cost at most $\exp(O(K\log K))$, while the tilt telescopes as above. Thus
\[
 \log a_n\ge n\log\mu-O(K\log K).
\]
By \eqref{eq:LK} and $R=O(h_K)$, one has
$K=\Theta((n/\log n)^{1/3})$, hence $K\log K=\Theta(\Phi(n))$. This proves the lower half of \eqref{eq:main} for all large integers, not merely a subsequence.

\section{A finite height radius gain}
\subsection{A fixed gap derivative bound}
For fixed $t=t_\st$ and variable $x$, put $b=1-x$, $\beta=xt_\st/b$ and $\eta=x^2t_\st/b^2$. The coefficient formula \eqref{eq:lagrange} implies
\begin{equation}\label{eq:fixedgap}
 Q_q(x,t_\st)=1+x\sum_{j=1}^q W_j(x),\qquad
 W_j(x)=\left(\frac{1+\beta}{b}\right)^j
 \frac{N_{j-1}(\eta)}{(1-\eta)^{2j-1}},
\end{equation}
where $N_0=1$ and, for $j\ge2$,
\[
 N_{j-1}(v)=\frac1{j-1}\sum_{l=0}^{j-2}
 \binom{j-1}{l}\binom{j-1}{l+1}v^l.
\]
Here is a formal proof, with no analytic continuation needed. If $F_j(v)=\sum_{k\ge0}\mathcal C(j,k)v^k$, its coefficient ratio gives
\[
 v(1-v)F_j''+[2-(2j+2)v]F_j'-j(j+1)F_j=0.
\]
Substitute $F_j=(1-v)^{1-2j}N$. The equation becomes
\[
 v(1-v)N''+[2-(4-2j)v]N'-(2-j)(1-j)N=0.
\]
With $N(0)=1$, its unique formal coefficients satisfy
\[
 (l+1)(l+2)n_{l+1}=(l+2-j)(l+1-j)n_l.
\]
They are exactly the displayed coefficients of $N_{j-1}$, and terminate after $j-2$. For $j=1$, $\mathcal C(1,k)=1$ directly gives $F_1=(1-v)^{-1}$.

Since $\eta(\rho)=0.12737466\ldots<1$, choose a fixed interval $[\rho,\rho+\epsilon_0]$ on which $b>0$ and $0<\eta<1$. The positive polynomial $N_{j-1}$ has degree at most $j$, so
\[
 0\le\frac{\eta N_{j-1}'(\eta)}{N_{j-1}(\eta)}\le j.
\]
Logarithmic differentiation of \eqref{eq:fixedgap} gives a constant $C_0$, independent of $j$, such that $0\le(\log W_j)'\le C_0j$ throughout this interval. Increasing a constant $C_1$ therefore gives
\begin{equation}\label{eq:derivative}
 Q_q(x,t_\st)\le Q_q(\rho,t_\st)
       \exp(C_1(q+1)(x-\rho))
\end{equation}
for all $q\ge1$ and $x$ in the interval.

\subsection{The critical jump tail}
At $x=\rho,t=t_\st$, let $b=b_\st$ and $s=b^2-\rho^2t_\st>0$. The discriminant of \eqref{eq:Qquad} as a quadratic in $Q$ factors, as a polynomial in $z$, as
\[
 \Delta(z)=s^2(1-z/z_\st)(1-z/z_2),\qquad
 z_2=\frac{s^2}{(b+\rho t_\st)^2z_\st}.
\]
The exact certificate verifies $0<z_\st<z_2<1$, with $z_2=0.8817230195\ldots$. The branch with $Q(0)=0$ is
\begin{equation}\label{eq:sqrt}
 Q(\rho;z,t_\st)=
 \frac{s-(b-\rho t_\st)z-s\sqrt{1-z/z_\st}\sqrt{1-z/z_2}}
      {2\rho t_\st(1-z)}.
\end{equation}
The rational nonsingular part is analytic in $|z|<1$, and the multiplier of $\sqrt{1-z/z_\st}$ is analytic in $|z|<z_2$. Consequently, for some fixed $C_2$,
\begin{equation}\label{eq:qtail}
 0\le Q_q(\rho,t_\st)z_\st^q\le C_2q^{-3/2}\qquad(q\ge1).
\end{equation}
For an elementary coefficient justification, rescale $z=z_\st w$. The coefficients of $\sqrt{1-w}$ have absolute value at most $C(m+1)^{-3/2}$. The analytic multiplier has coefficients $O(R^{-m})$ for some $R>1$. In the convolution, indices at most $q/2$ yield $O(q^{-3/2})$ by geometric summation, and the other indices yield an exponentially small bound. The rational part is exponentially smaller as well. This proves \eqref{eq:qtail} without a transfer theorem.

\subsection{The cutoff estimate}
Let $a_n^{[\le H]}$ denote the coefficients in \eqref{eq:walk} restricted to macro-boundary heights in $\{1,\ldots,H\}$. This is a coefficientwise positive subclass of the formal walk expansion. Set
\[
 \delta_H=\epsilon\frac{\log(H+1)}{H+1},\qquad x_H=\rho+\delta_H,
\]
with a sufficiently small fixed $\epsilon>0$. Write $A_q=Q_q(\rho,t_\st)z_\st^q$. By \eqref{eq:derivative}, the weighted row sum of the cutoff at $x_H$ is at most
\begin{equation}\label{eq:rowcut}
 \frac{x_H}{1-x_H}t_\st e^{C_1\delta_H}
       \sum_{q\le H}A_qe^{C_1q\delta_H}.
\end{equation}
Its unperturbed infinite mass is at most $m_\st<1$. Moreover, \eqref{eq:qtail} implies
\begin{equation}\label{eq:tailperturb}
 \sum_{q\le H}A_q(e^{C_1q\delta_H}-1)\longrightarrow0.
\end{equation}
Indeed split at $q=1/\delta_H$. The lower part is $O(\sqrt{\delta_H})$, using $e^{C_1q\delta_H}-1=O(q\delta_H)$ there. The upper part is bounded by
\[
 C e^{C_1H\delta_H}\sqrt{\delta_H}
 =O\!\left(H^{C_1\epsilon-1/2}\sqrt{\log H}\right)=o(1)
\]
if $C_1\epsilon<1/2$. The prefactor in \eqref{eq:rowcut} tends to its critical value. Thus the row mass is uniformly at most a fixed $\bar m<1$ for large $H$. The same weighted geometric bound as in Lemma~\ref{lem:criticalbound}, now in the finite cutoff, gives a uniform upper bound on its generating function at $x_H$. Positive coefficient extraction yields
\begin{equation}\label{eq:cutbound}
 a_n^{[\le H]}\le Cx_H^{-n}
 \le C\mu^n\exp\!\left(-c\frac{n\log(H+1)}{H+1}\right).
\end{equation}
All constants are independent of $H,n$ once $H$ is sufficiently large.

\section{The first hit bound and the upper estimate}
Keep $W=W_\st$ fixed and write the right side of \eqref{eq:W} as
\[
 F(x,z,t;W)=\frac{z}{1-x}
 \frac{(1+xt/(1-x))(1+W)}{1-x^2t(1+W)/(1-x)^2}.
\]
Let $f(x,\theta)=F(x,z_\st e^{-\theta},t_\st e^\theta;W_\st)$. By \eqref{eq:margins}, $f(\rho,0)=W_\st$ and $f_\theta(\rho,0)=0$. Specifically, with $u=\eta_\st(1+W_\st)$,
\[
 \frac{f_\theta}{f}(\rho,0)
 =-\frac1{1+\beta_\st}+\frac{u}{1-u}=0.
\]
Also $f_x(\rho,0)>0$, since the positive factors increase with $x$ and the denominator remains positive. Choose a fixed sufficiently large $D>0$ and put
\[
 x_\theta=\rho e^{-D\theta^2},\qquad
 t_\theta=t_\st e^\theta,\qquad z_\theta=1/t_\theta.
\]
Taylor's theorem gives
\[
 f(x_\theta,\theta)=W_\st+
 \left(\tfrac12 f_{\theta\theta}(\rho,0)-\rho D f_x(\rho,0)\right)\theta^2
 +O(\theta^3).
\]
Make the bracket negative. For $0\le\theta\le\theta_0$ with sufficiently small $\theta_0>0$, the positive fixed-point map then has the supersolution $W_\st$ and positive denominator. Iteration from zero gives
\[
 Q(x_\theta;z_\theta,t_\theta)
 \le\frac{z_\theta+x_\theta W_\st}{1-z_\theta}.
\]
Its unrestricted weighted row mass satisfies
\[
 m_\theta=\frac{x_\theta}{1-x_\theta}t_\theta
 Q(x_\theta;z_\theta,t_\theta)
 \le \frac{x_\theta}{1-x_\theta}t_\theta
       \frac{z_\theta+x_\theta W_\st}{1-z_\theta}.
\]
The explicit supersolution bound on the right is continuous at zero, where it equals $m_\st<1$. Reducing $\theta_0$ makes all such masses at most one fixed $\bar m<1$. The sum of the weights of all finite macro-prefixes starting at height $1$, including the factor $t_\theta^h$ at their final height but no terminating weight, is therefore at most
\begin{equation}\label{eq:prefix}
 t_\theta\sum_{k\ge0}\bar m^k\le C.
\end{equation}

Let $p_{m,h}$ count the first-hit macro-prefixes of length $m$ whose first height above $H$ is $h>H$. The remaining unrestricted suffix has generating function $e_h$. By \eqref{eq:ehcritical}, for the complementary coefficients at fixed word length $n$,
\[
 a_n^{[>H]}\rho^n
 \le C\rho\sum_{m\le n-1,\,h>H}p_{m,h}\rho^m t_\st^h.
\]
Here $\rho$ accounts for the original leading letter, and the use of $e_h(\rho)$ only enlarges the permitted remaining lengths. For $0<\theta\le\theta_0$,
\[
 \rho^m\le x_\theta^m e^{D\theta^2n},\qquad
 t_\st^h\le t_\theta^h e^{-\theta H}.
\]
First-hit prefixes are a subset of all prefixes, so \eqref{eq:prefix} implies
\[
 a_n^{[>H]}\le C\mu^n\exp(D\theta^2n-\theta H).
\]
When $H/(2Dn)\le\theta_0$, use $\theta=H/(2Dn)$ to obtain
\begin{equation}\label{eq:hit}
 a_n^{[>H]}\le C\mu^n\exp\!\left(-\frac{H^2}{4Dn}\right).
\end{equation}
Continuation at the critical tilt in \eqref{eq:ehcritical} is what cancels the final height weight. No constraint on the final height of the entire walk is imposed.

Take $H=\lfloor n^{2/3}(\log n)^{1/3}\rfloor$. It is $o(n)$, so \eqref{eq:hit} applies for large $n$. Both exponents in \eqref{eq:cutbound} and \eqref{eq:hit} are at least a positive constant times $\Phi(n)$. Adding the two positive subclasses proves the upper half of \eqref{eq:main}. Together with the lower bound, this completes Theorem~\ref{thm:main}.

\section{Inverse order and non D finiteness}
The main theorem gives $c\Phi(n)-O(1)\le n\log\mu-\log a_n\le C\Phi(n)$. If a pure stretched equivalent in Corollary~\ref{cor:consequences}(1) held, its deficit would be $dn^\gamma-g\log n-\log C_0+o(1)$. This grows faster than $\Phi(n)$ for $\gamma>1/3$ and more slowly for $\gamma=1/3$, contradicting the two-sided bounds.

The sequence is nondecreasing: append a copy of the last letter to an avoiding ascent sequence. The ascent condition persists. A new 120 occurrence using the appended letter would already occur using its previous copy. This extension is injective. Hence $N(Y)$ is well defined for large $Y$.

Let $L=\log Y$ and $\lambda=\log\mu$. Applying the upper bound at $n=N(Y)$ first gives $N(Y)\ge L/\lambda-O(1)$ and then
\[
 \lambda N(Y)-L\ge c\Phi(N(Y))-O(1)\ge c'\Phi(L).
\]
Conversely set $n=\lceil L/\lambda+K\Phi(L)\rceil$. Since $\Phi(n)/\Phi(L)\to\lambda^{-1/3}$, any sufficiently large fixed $K$ makes $\lambda n-C\Phi(n)\ge L$. The coefficient lower bound implies $a_n\ge Y$, hence $N(Y)\le n$. This proves the stated inverse order. It is a positive correction of a known order, with no asserted leading coefficient or exact rounding rule.

For non-D-finiteness, the stretched upper bound gives
\[
 \sum_{n\ge0}n^k a_n\rho^n<\infty\qquad(k=0,1,2,\ldots).
\]
Thus $A$ and all its derivatives extend continuously along the real radius to $\rho$. Suppose $A$ were D-finite over $\Q(x)$. Its coefficients are exponentially bounded nonnegative integers, their algebraic conjugates have the same bound, and all common denominators are $1$. It would therefore be a Siegel G-function. The G-function regularity theorem gives a convergent finite local sum of terms
\[
 (x-\rho)^\alpha(\log(x-\rho))^\beta H(x-\rho),
\]
with rational $\alpha$, nonnegative integer $\beta$ and analytic $H$, on a slit neighborhood of $\rho$; see \cite[Definition 2.2 and Theorem 2.1(c)]{GB}.

Such a convergent finite power-log expansion can be smooth to every order along a radius only if it is analytic at the endpoint. To see this, combine terms with the same rational exponents modulo integers and the same logarithmic powers. The possible exponents form a discrete set bounded below. A least surviving negative integer power is unbounded; a least noninteger power fails to have a bounded derivative at some finite order; and a least nonzero logarithmic term fails similarly. The absence of all these terms leaves a convergent ordinary power series. Consequently $A$ would be analytic at $\rho$, contrary to Pringsheim's theorem for a nonnegative series of finite radius. Thus $A$ is not D-finite over $\Q(x)$.

Finally, a complex differential relation of bounded order and polynomial degree for a rational series produces a homogeneous linear system with rational entries and finitely many unknowns. Its rank is attained by finitely many rows. If it has a nonzero complex solution, it has a nonzero rational solution. Hence complex D-finiteness would imply rational D-finiteness here, proving Corollary~\ref{cor:consequences}(3). This argument uses G-function arithmetic regularity; it does not assert that all D-finite functions have regular singularities.

\section{Reproduction and remaining questions}
The exact certificate checks the algebraic identities and rational sign margins. Independent programs compare literal words, normalized states, gap operators and binomial height walks, including the OEIS terms through $n=20$. Additional checks cover fixed-gap formulas and staircase legality, lengths and tilt cancellation. These diagnostics supplement the proofs above.

The package includes editable source, a build script, mathematical notes, a separate audit and content hashes. The checks require Python with SymPy and mpmath. The PDF build uses standard LaTeX with AMS packages, geometry, hyperref, microtype and Latin Modern fonts.

The convergence and possible limit of $\bigl(n\log\mu-\log a_n\bigr)/\Phi(n)$, and lower-order corrections, remain open here. There is no full transseries or all-orders inverse expansion. Although the $n^{3/8}$ estimate is ruled out asymptotically, the theorem gives no explicit crossover threshold for finite-data fits.

\begin{thebibliography}{9}
\small
\bibitem{OEIS} OEIS Foundation, \emph{A202061: Number of ascent sequences avoiding the pattern 120}. Inspected 2 October 2026. \url{https://oeis.org/A202061}.
\bibitem{Conway} A. R. Conway, M. Conway, A. Elvey Price and A. J. Guttmann, \emph{Pattern-Avoiding Ascent Sequences of Length 3}, Electronic Journal of Combinatorics \textbf{29}(4) (2022), P4.25. \url{https://www.combinatorics.org/ojs/index.php/eljc/article/download/v29i4p25/pdf/}.
\bibitem{GB} S. Garoufalidis and J. Bellissard, \emph{Algebraic G-functions associated to matrices over a group-ring}, version dated 19 December 2013, Definition 2.2 and Theorem 2.1(c). \url{https://people.mpim-bonn.mpg.de/stavros/publications/algebraicGfunctions.pdf}.
\end{thebibliography}
\end{document}
